% ==================================================================
%  SETUP + DECOMPOSITION (pp. 2--3): Pillar 1
% ==================================================================

\section{Setup and the Four-Term Decomposition}
\label{sec:decomposition}

\subsection{Problem Setup}
\label{sec:setup}

A \emph{budgeted transductive task} consists of an instance space~$\cX$,
a solution space~$\cY$ with an external verifier
$v\colon \cX \times \cY \to \{0,1\}$, and a resource budget~$B$.
The agent receives $x \in \cX$ and must produce~$y$ with $v(x,y)=1$
within budget~$B$.  Canonical examples include code repair verified against
test suites and formal proofs checked by a proof assistant.

A \emph{design} is a tuple $d = (M, D, \kappa)$ specifying a model~$M$,
a retry depth~$D$, and a per-step cost~$\kappa(M)$.  The \emph{design
menu} $\cD = \{d_1,\ldots,d_k\}$ enumerates the candidates.  Our goal is
to select the design in~$\cD$ that maximizes expected performance
(verified-solution probability) subject to a cost--quality tradeoff.

\begin{assumption}[Packed latent-mode family]
\label{ass:packed}
The task distribution admits a latent partition into $\Mn$ modes
with distribution $\pi = (\pi_1,\ldots,\pi_{\Mn})$ such that:
(i)~each instance has a unique dominant mode;
(ii)~conditional on mode~$s$, the number of attempts to a verified
solution under model~$M$ follows
$T(s,M) \sim \mathrm{Geometric}(p(s,M))$;
(iii)~the mode distribution~$\pi$ does not shift with model choice.
\end{assumption}

The assumption holds when task difficulty clusters into discrete types
(e.g., repository-specific bug patterns in SWE-bench, or theorem
categories in formal mathematics) and each type admits a memoryless solve
probability.  Outside the packed family, the decomposition below holds with
a bounded residual of order
$O(\delta \log \Mn + \eta)$~\citep{beneventano2026framework}.

\subsection{The Four-Term Identity}
\label{sec:four-term}

\begin{theorem}[Four-term decomposition]
\label{thm:four-term}
Under Assumption~\ref{ass:packed}, the log-performance gap between designs
$(M_0, q_0)$ and $(M, q)$ decomposes as
\begin{equation}
\label{eq:four-term}
  \Delta \;=\;
  \underbrace{\log\frac{\kappa_0}{\kappa(M)}}_{\text{cost}}
  \;+\; \underbrace{\sum_s \pi_s \log\frac{p(s,M)}{p(s,M_0)}}_{\varphi\;\text{(competence)}}
  \;+\; \underbrace{\MI_\pi(S;\,D_M)}_{G\;\text{(info gain)}}
  \;-\; \underbrace{\E_D\!\bigl[\KL(\pi_D^M \| q_D)\bigr]}_{\mismatch\;\text{(mismatch)}}.
\end{equation}
Each term is independently estimable from pilot data.
\end{theorem}

\begin{proof}[Proof sketch]
Expand $\E[-\log T]$ under the geometric-trials model and rewrite the
routing allocation via the chain rule for KL divergence.  The cost term
separates because $\kappa$ does not depend on the task mode.  The
competence term collects the per-mode log-probability ratios.  The
remaining cross-terms decompose into mutual information (information
gain) and KL divergence (routing mismatch) by standard
identities~\citep{beneventano2026framework}.
Full proof in Appendix~\ref{app:four-term}.
\end{proof}

The \textbf{cost term} measures the per-step price difference.  The
\textbf{competence term}~$\varphi$ measures how much better model~$M$ is
at each mode, averaged over the task distribution.  The \textbf{information
gain}~$G$ measures the benefit of mode-aware routing---it is the value of
knowing which task type you face, in the sense of
Kelly~\citeyearpar{kelly1956new}.  The \textbf{mismatch} measures how far
actual routing deviates from the posterior-optimal allocation.

\subsection{The Crossover Formula}
\label{sec:crossover}

\begin{theorem}[Crossover depth]
\label{thm:crossover}
Consider a large model with per-attempt success probability~$p_l$ and
per-step cost~$\kappa_l$, and a small model with $p_s < p_l$ and
$\kappa_s < \kappa_l$.
On a single-mode task, the small model with $D$ retries achieves the
same expected verified-solution probability as the large model (single
attempt) when
\begin{equation}
\label{eq:crossover}
  D_{\mathrm{cross}} = \frac{\log(1 - p_l)}{\log(1 - p_s)}.
\end{equation}
The small model dominates whenever $D_{\mathrm{cross}} \cdot \kappa_s < \kappa_l$.
\end{theorem}

\begin{proof}
The success probability after $D$ retries is $1 - (1-p_s)^D$.
Setting this equal to~$p_l$ and solving for~$D$ gives~\eqref{eq:crossover}.
\end{proof}

As a worked example: if $p_l = 0.50$ and $p_s = 0.10$, then
$D_{\mathrm{cross}} = \log(0.50)/\log(0.90) \approx 6.6$.
If the large model costs $10\times$ more per step, the small model
with 7 retries is both cheaper and more accurate.

\begin{corollary}[Monotonicity]
\label{cor:monotonicity}
Under Assumption~\ref{ass:packed}, the optimal depth $D^*(M)$ is
non-increasing in the per-mode success probability: stronger models need
fewer retries.  For fixed depth, the expected performance is monotonically
increasing in $p(s,M)$ for every mode.
\end{corollary}

\begin{remark}[The Lagrangian non-connection]
\label{rem:lagrangian}
The four terms in~\eqref{eq:four-term} are components of the primal
objective evaluated at any feasible point, not dual variables of a
Lagrangian relaxation.  A natural conjecture that they are shadow prices
is precisely wrong: the KKT conditions yield a single scalar multiplier,
not four (Appendix~\ref{app:lagrangian}).
\end{remark}
